\chapter{Einleitung}
\label{chap:01_einleitung}
\pagenumbering{arabic}
   
In dieser Arbeit wird ein System vorgestellt, mit dem Handbewegungen aufgenommen werden können. Das System wurde dafür in zwei Einzelsysteme aufgeteilt, zum einen für die Bestimmung der absoluten Position der Hand und zum anderen die Bestimmung der Fingerhaltung relativ zum Handrücken. 

In diesem Kapitel wird das Thema motiviert und diese Arbeit in einen aktuellen Kontext gesetzt. Abschließend wird die Gliederung dieser Arbeit vorgestellt.

   \section{Motivation}
   \label{sec:01_motivation}
   Die Handbewegungsanalyse spielt in verschiedenen Bereichen eine Rolle. In der Medizin wird diese erforscht um Parkinson zu behandeln ref. In der Unterhaltungsindustrie kommen bei der Animation von Bewegungen zum Einsatz und in VR/AR-Anwendungen können diese als Schnittstelle zwischen Mensch und Maschine verwendet werden (HMI). 
   In der Forschung werden unterschiedliche Ansätze erforscht. Für das Bewegungstracking werden unterschiedliche Sensoren erforscht. Die Bewegungserfassung erfolgt auf Basis von Optischen-, Stretch- oder magnetischen Sensorsystemen. 
      
   \section{Einordnung der Arbeit}
   \label{sec:01_einordnung}

      
   \section{Aufbau der Arbeit}
   \label{sec:01_aufbau}
